\documentclass[11pt]{article}
\usepackage[margin=1in]{geometry}
\usepackage{amsmath,amssymb,amsthm,hyperref}
\newtheorem{theorem}{Theorem}
\newcommand{\conv}{\operatorname{conv}}
\newcommand{\W}{\mathcal W}
\title{Simplicity of even signed Wachs permutation polytopes}
\author{Principia Math In Progress}
\date{27 September 2026}
\begin{document}
\maketitle
\begin{abstract}
The convex hull of the signed Wachs permutations in $B_{2m}$ is
linearly isomorphic to a signed permutahedron times an $m$-cube.
An explicit cone computation proves simplicity in dimension $2m$ for
every $m\geq1$, settling Conjecture 6.4 of Brenti and Sentinelli.
\end{abstract}

\section*{Task 1: the signed conjecture}
A signed permutation $w\in B_n$ is a bijection of
$[\pm n]=\{-n,\ldots,-1,1,\ldots,n\}$ satisfying $w(-i)=-w(i)$.
Its coordinate vector is $(w(1),\ldots,w(n))\in\mathbb R^n$.
For $i\in[2m]$ put $i^*=i+1$ if $i$ is odd and $i^*=i-1$ if $i$ is even.
The source defines
\[
 \W(B_{2m})=\{w\in B_{2m}:|w^{-1}(i)-w^{-1}(i^*)|\leq1
                  \text{ for every }i\in[2m-1]\}.
\]
Let $Q_m=\conv\W(B_{2m})$. Conjecture 6.4 of \cite{BS} says that
$Q_m$ is simple for every positive integer $m$.

\begin{theorem}\label{main}
Put $a=(3,7,\ldots,4m-1)$ and let
\[
 \Pi^B(a)=\conv\{(\epsilon_1a_{\sigma(1)},\ldots,
             \epsilon_ma_{\sigma(m)}):\sigma\in S_m,\ \epsilon_i\in\{-1,1\}\}.
\]
Then $Q_m$ is linearly isomorphic to $\Pi^B(a)\times[-1,1]^m$.
It is a simple $2m$-dimensional polytope with $4^m m!$ vertices
and $m4^m m!$ edges.
\end{theorem}

\section*{Signed blocks and the product map}
For each $j$, the signed positions $w^{-1}(2j-1)$ and $w^{-1}(2j)$
are distinct and have distance one. They must have the same sign:
nonzero integers of opposite signs have distance at least two.
Therefore, in the positive-position coordinate vector, the values with
absolute values $2j-1,2j$ occur consecutively and with the same sign.
These disjoint absolute-position pairs cover $[2m]$; position 1 must
be paired with 2, and induction gives precisely
$\{1,2\},\{3,4\},\ldots,\{2m-1,2m\}$.
Conversely, every arrangement of the $m$ value pairs, with an
independent common sign and an independent orientation for each pair,
satisfies the source's signed inverse-position condition.

Define a linear isomorphism $T:\mathbb R^{2m}\to\mathbb R^m\times\mathbb R^m$
by
\begin{equation}\label{map}
 y_i=x_{2i-1}+x_{2i},\qquad z_i=x_{2i}-x_{2i-1}.
\end{equation}
Its inverse is $x_{2i-1}=(y_i-z_i)/2$, $x_{2i}=(y_i+z_i)/2$.
A block with absolute values $2j-1,2j$ has sum $\pm(4j-1)$ and
difference $\pm1$. The signs of its sum and difference can be chosen
independently, because its common sign and orientation are independent.
Thus $T(\W(B_{2m}))$ is exactly the Cartesian product of the signed
permutation vectors of $a$ and $\{-1,1\}^m$.
For finite sets $A,B$, the identity
$\conv(A\times B)=\conv(A)\times\conv(B)$ follows by multiplying
the weights of convex combinations. Consequently,
\begin{equation}\label{product}
 T(Q_m)=\Pi^B(a)\times[-1,1]^m.
\end{equation}

\section*{A direct cone proof for the signed factor}
Every signed permutation vector $u$ of $a$ has the same squared norm
as $a$. Hence
\[
 a\cdot u=\|a\|^2-\tfrac12\|u-a\|^2,
\]
so the functional $u\mapsto a\cdot u$ uniquely exposes $a$ as a vertex.
Signed coordinate permutations show that every generating vector is a
vertex and that the signed symmetry group acts transitively on them.

For such a vector $u$, define
\[
 c_0=\sum_{i=1}^m(a_i-u_i),\qquad
 c_k=\sum_{i=k+1}^m(a_i-u_i)\quad(1\leq k<m).
\]
All these numbers are nonnegative. Indeed, the sum of any $r$ entries
of a signed permutation of the positive increasing vector $a$ is at most
the sum of the $r$ largest entries of $a$. Telescoping now gives
\begin{equation}\label{cone}
 u-a=c_0(-e_1)+\sum_{k=1}^{m-1}c_k(e_k-e_{k+1}).
\end{equation}
The ray $\mathbb R_{\geq0}(-e_1)$ is attained by changing the sign of
$a_1=3$, giving the vertex difference $-6e_1$.
For each $k$, interchanging $a_k,a_{k+1}$ gives the difference
$4(e_k-e_{k+1})$. Therefore the feasible-direction cone at $a$ is
exactly the nonnegative span of these $m$ linearly independent vectors.
In particular the signed factor has dimension $m$.

At $(a,\mathbf1)$ in \eqref{product}, the cone is generated by
\[
 (-6e_1,0),\quad (4(e_k-e_{k+1}),0)\ (1\leq k<m),\quad
 (0,-2e_i)\ (1\leq i\leq m).
\]
Equation \eqref{cone} gives coefficients $c_0/6,c_k/4$ for the first
factor, and the cube coefficients are $(1-z_i)/2$. Every generator is
an attained vertex difference. The displayed $2m$ generators are linearly
independent, and therefore constitute exactly the extreme rays of the
cone. Extreme rays of a vertex cone correspond to edges through that
vertex, by exposing the ray or the edge with a linear functional.
Thus the product has dimension $2m$ and exactly $2m$ edges at this vertex.
Its signed coordinate symmetries and cube symmetries act transitively
on all vertices, proving simplicity everywhere. The signed factor has
$2^m m!$ vertices, so the product has $4^m m!$ vertices. Counting
vertex--edge incidences gives $m4^m m!$ edges. When $m=1$ this is a
parallelogram. Theorem \ref{main} follows.

\section*{Scope and verification}
This proves the complete even signed conjecture. It is distinct from
the unsigned Conjecture 6.3: the unsigned hull is a facet of $Q_m$,
exposed by maximizing the sum of coordinates. Simplicity of that facet
alone does not prove simplicity of $Q_m$. The additional signed
position argument and the full-dimensional cone proof above supply
the signed resolution without assuming the unsigned theorem.

The exact verifier exhausts $B_{2m}$ for $1\leq m\leq3$ using the
original signed inverse-position condition. It checks the map, inverse
and cone decompositions on all product vertices through $m=4$.
These are finite consistency checks; the proof for every $m$ is above.
Review is by the authoring agent, without independent peer review or
proof-assistant certification.

\section*{Proposed follow-up tasks}
\begin{description}
\item[Task 2.] Determine all face numbers of $Q_m$ in closed form from
the signed ordered-partition description of $\Pi^B(a)$ and the cube factor.
\item[Task 3.] Compute the Ehrhart polynomial of $Q_m$ in
$\mathbb Z^{2m}$, explicitly describing the parity constraints imposed
by the inverse of \eqref{map}.
\item[Task 4.] Determine the full affine symmetry group of $Q_m$ for
every $m$, including any exceptional symmetries in low dimensions.
\item[Task 5.] Classify the odd integers $n$ for which
$\conv\W(B_n)$ is simple, and give explicit nonsimple vertices for
every other odd $n$.
\end{description}

\begin{thebibliography}{9}
\bibitem{BS} F. Brenti and P. Sentinelli,
\emph{Wachs permutations, Bruhat order and weak order},
European Journal of Combinatorics \textbf{119} (2024), 103804,
\href{https://doi.org/10.1016/j.ejc.2023.103804}{doi:10.1016/j.ejc.2023.103804}.
\href{https://arxiv.org/abs/2212.04932}{arXiv:2212.04932v1};
definition on p.~6, signed block bijection on p.~19, Conjecture 6.4
on pp.~36--37. The accompanying audit records the fully read
institutional postprint and auxiliary author-preprint versions.
\end{thebibliography}
\end{document}
